% !TEX root = ../thesis.tex
\chapter*{TÓM TẮT KHÓA LUẬN}
\addcontentsline{toc}{chapter}{TÓM TẮT KHÓA LUẬN}

Sự bùng nổ của thông tin trực tuyến tại Việt Nam đã đặt ra hai nhu cầu song song: tóm tắt
nhanh nội dung bài báo cho người đọc, và kiểm chứng độ tin cậy cho
người đọc muốn xác thực. Các hệ thống tóm tắt dựa trên một Mô hình ngôn ngữ lớn
(LLM) đơn lẻ hiện nay vẫn cho ra kết quả không ổn định, trong khi kỹ thuật
\textbf{Mixture-of-Agents} (MoA) --- kết hợp nhiều LLM thông qua kiến trúc tác nhân
đề xuất--tổng hợp --- chưa từng được nghiên cứu và ứng dụng cho bài toán tóm tắt tin
tức tiếng Việt. Bên cạnh đó, các hệ số đo lường được sử dụng phổ biến nhất trong
nghiên cứu tiếng Việt (ROUGE, BLEU, BERTScore) đo độ trùng lặp $n$-gram hoặc độ
tương đồng nhúng từ với bài báo gốc, vốn là chỉ số kém hiệu quả để phản ánh chất
lượng tóm tắt khi mô hình được phép diễn giải lại và rút gọn ý tưởng nhằm phục vụ mục đích tóm tắt.

Khóa luận này thực hiện \textbf{ba đóng góp}. \textit{Thứ nhất}, khóa luận giới thiệu
\textbf{Fiber} --- một tiện ích mở rộng trình duyệt và hệ thống backend Next.js có khả
năng tóm tắt và kiểm chứng tin tức tiếng Việt theo thời gian thực, hỗ trợ các LLM
từ OpenAI, Anthropic, Google và cả HuggingFace. \textit{Thứ hai}, khóa luận áp dụng
pipeline Mixture-of-Agents của Wang et al.\ (2024) cho bài toán tóm tắt tiếng Việt có
neo ngữ cảnh gốc, bằng cách đưa bài viết nguồn trực tiếp vào ngữ cảnh của tác nhân
tổng hợp như một \textit{kết nối phần dư} (residual connection) nhằm đảm bảo bản tóm
tắt cuối cùng trung thực với bài gốc. \textit{Thứ ba}, và quan trọng nhất, khóa luận
đề xuất một \textbf{khung đánh giá ba trục}: Trục A (giữ lại nội dung: ROUGE, BLEU,
BERTScore), Trục B (chất lượng và sở thích tự động: đánh giá rubric, xếp hạng từng
cặp bởi LLM-giám khảo, phân loại tính xác thực), và Trục C (đánh giá con người: xếp
hạng mù $K$-chiều có tính chỉ số đồng thuận giữa các người đánh giá), đồng thời dùng
khung đó để chỉ ra vì sao các hệ số overlap riêng lẻ không thể đánh giá đầy đủ chất
lượng tóm tắt.

Trên tập dữ liệu cân bằng chủ đề gồm 154 bài báo từ \textit{tienphong.vn}, pipeline
MoA vượt trội so với mô hình đơn lẻ mạnh nhất (lựa chọn bản thảo tốt nhất) trong
93 trên 135 phán quyết quyết định ($68,9\%$, $p < 0,0001$); đồng thời
cũng vượt trội so với từng mô hình đề xuất riêng lẻ. Điểm chất lượng rubric và các
hệ số overlap đều đồng thuận với kết quả này. Một nghiên cứu đánh giá con người toàn
diện với 17 người đánh giá trên tất cả 48 tác vụ (192 xếp hạng) cho thấy chênh lệch lớn giữa sở thích của LLM-giám khảo và sở thích của con người. Khoảng cách này,
cùng với sự tương tác phụ thuộc chủ đề (con người ưa bản thảo ngắn hơn với nội dung
hành chính, nhưng ưa bản tóm tắt dung hợp với nội dung giải trí), chính là phát hiện
phương pháp luận mà khung đánh giá được xây dựng để khẳng định.

\vspace{0.4em}
\textbf{Từ khóa:} Tóm tắt tin tức tiếng Việt, kiểm chứng thông tin, mô hình ngôn ngữ
lớn, Mixture-of-Agents, tiện ích mở rộng trình duyệt, khung đánh giá
ba trục.
